\chapter{Observabilidad, información de Fisher e incertidumbre}
\label{ch:observability}

Los capítulos precedentes construyeron el modelo directo, desde la escena hasta la
medición. Este capítulo estudia su dirección inversa: bajo qué condiciones una variación
del estado humano produce cambios distinguibles en los observables. La respuesta se
formula localmente mediante sensibilidad, información estadística e incertidumbre, y
servirá después para evaluar tanto los algoritmos como el diseño de la red.

\section{Problema inverso, identificabilidad y observabilidad local}

El modelo directo lleva un estado $\vect q$ a una distribución de observaciones $p(\vect y\mid\vect q)$. La inversión pregunta qué estados son compatibles con una medición. La \emph{identificabilidad global} exige que estados distintos no induzcan la misma distribución; la \emph{identificabilidad local} estudia esa unicidad en una vecindad; y la \emph{estimabilidad práctica} incorpora ruido, número finito de muestras y error del modelo \citep{KaipioSomersalo2005,Kay1993}. En este manuscrito, \emph{observabilidad} denota información estadística local sobre variables físicas, no el criterio clásico de observabilidad de un sistema dinámico lineal.

Sea
\begin{equation}
\vect y=\vect\mu(\vect q)+\vect w,
\qquad
\vect w\sim\mathcal{CN}(\vect 0,\mat\Sigma_w).
\label{eq:local_observation}
\end{equation}
La matriz Jacobiana
\begin{equation}
\mat J_q=\frac{\partial\vect\mu}{\partial\vect q^T}
\label{eq:state_jacobian}
\end{equation}
describe sensibilidad, pero su norma no considera ruido ni correlación. Para parámetros reales, ruido complejo propio, covarianza conocida e independiente de $\vect q$, la \ac{FIM} es
\begin{equation}
\mat F_q(\vect q)
=2\Re\left\{
\mat J_q^H\mat\Sigma_w^{-1}\mat J_q
\right\}.
\label{eq:fim}
\end{equation}
Si la covarianza depende del estado aparecen términos adicionales; si el ruido no es gaussiano, la FIM se deriva de la función de puntuación y no de la \cref{eq:fim}. Bajo condiciones regulares, para parámetros identificables y estimadores localmente
insesgados, la cota de Cramér--Rao es
$\operatorname{Cov}(\widehat{\vect q})\succeq\mat F_q^{-1}$
\citep{Kay1993}. Si la información es singular, se declara un estimando local
$U^T\vect q$ en un subespacio identificable de base ortonormal $U$ y se aplica
la cota correspondiente a esa parametrización regular. La pseudoinversa no
asigna incertidumbre nula a una dirección invisible y no garantiza la existencia
de un estimador insesgado del vector completo. La cota tampoco equivale a un
intervalo posterior ni incorpora automáticamente discrepancia del gemelo.

\section{Perturbaciones e información equivalente}

Sea $\vect\zeta$ un vector de perturbaciones que contiene, por ejemplo, la fase común, la ganancia o la incertidumbre geométrica. La FIM conjunta se particiona como
\begin{equation}
\mat F(\vect q,\vect\zeta)=
\begin{bmatrix}
\mat F_{qq}&\mat F_{q\zeta}\\
\mat F_{\zeta q}&\mat F_{\zeta\zeta}
\end{bmatrix}.
\end{equation}
Cuando $\vect\zeta$ debe estimarse junto con $\vect q$, la información equivalente para el estado de interés es el complemento de Schur
\begin{equation}
\mat F_q^{\rm eq}=
\mat F_{qq}-
\mat F_{q\zeta}\mat F_{\zeta\zeta}^{\dagger}
\mat F_{\zeta q}.
\label{eq:equivalent_fim}
\end{equation}
Esta expresión formaliza por qué un canal sensible al movimiento puede ser poco informativo si el mismo cambio puede explicarse mediante CFO, fase desconocida o movimiento de otro sujeto.

\section{Escala, rango y criterios espectrales}

Los autovalores de una FIM dependen de las unidades de $\vect q$. Antes de comparar posición en metros, orientación en radianes y desplazamiento torácico en milímetros, se define
\begin{equation}
\bar{\vect q}=\mat D_q^{-1}(\vect q-\vect q_0),
\qquad
\mat F_{\bar q}=\mat D_q^T\mat F_q^{\rm eq}\mat D_q,
\label{eq:fim_normalization}
\end{equation}
con escalas físicas declaradas en $\mat D_q$. Para un conjunto fijo de variables
y $\mat F_{\bar q}\succ0$, se considerarán
\begin{align}
\lambda_{\min}(\mat F_{\bar q})
&\quad\text{(E-optimality)},\\
\log\det(\mat F_{\bar q})
&\quad\text{(D-optimality)},\\
\tr(\mat F_{\bar q}^{-1})
&\quad\text{(A-optimality)}.
\end{align}
Los criterios se comparan sobre el mismo subespacio objetivo; un diseño singular
se declara no admisible para recuperar todas sus variables. Añadir una matriz
de información previa define un criterio regularizado diferente y exige declarar
ese prior. En particular, minimizar la traza de una pseudoinversa sin controlar
el rango puede favorecer una configuración que pierde una variable.

El rango identifica el subespacio localmente visible; un valor propio pequeño indica una dirección mal condicionada, no necesariamente una variable individual invisible. Estos criterios provienen del diseño óptimo de experimentos y se reutilizarán para determinar la ubicación de los nodos \citep{Pukelsheim2006}.

\section{Respiración como ejemplo conductor}

Para $q=x_{\rm chest}$, y bajo las condiciones de la \cref{eq:fim}, el índice escalar
\begin{equation}
o_{\rm resp}(\vect x)=
2\Re\left\{
\left(\frac{\partial\vect\mu}{\partial x_{\rm chest}}\right)^H
\mat\Sigma_w^{-1}
\left(\frac{\partial\vect\mu}{\partial x_{\rm chest}}\right)
\right\}
\label{eq:resp_observability}
\end{equation}
define un mapa local de información del desplazamiento torácico con las
perturbaciones restantes conocidas. Si son desconocidas se utiliza la
información equivalente de la \cref{eq:equivalent_fim}. Su derivada hereda la geometría biestática de la \cref{eq:bistatic_sensitivity}; por ello el mapa puede diferir radicalmente de RSSI. La tesis validará si predice zonas favorables y puntos ciegos mejor que métricas de comunicación.

Presencia requiere un tratamiento diferente porque es una hipótesis discreta. Para $\mathcal H_0$ (ausencia) y $\mathcal H_1$ (presencia) se utilizarán razón de verosimilitudes, curvas ROC y separaciones distribucionales como KL o Chernoff, no una FIM ordinaria salvo que se introduzca y justifique una relajación continua.

\section{Redes distribuidas y dependencia estadística}

Si $N_{\rm link}$ enlaces son condicionalmente independientes dado $\vect q$, la información se suma:
\begin{equation}
\mat F_{\rm net}(\vect q)
=\sum_{l=1}^{N_{\rm link}}\mat F_{q,l}(\vect q).
\label{eq:fim_network}
\end{equation}
Con osciladores, interferencia o errores ambientales compartidos, la hipótesis falla y debe utilizarse la covarianza conjunta en la \cref{eq:fim}. Esta precaución distingue diversidad nominal de información realmente independiente.
Si los enlaces comparten una incógnita instrumental, se construye primero
la información conjunta y después se elimina esa incógnita; sumar
complementos de Schur independientes describe en general otro modelo.

\section{Observabilidad física y disponibilidad temporal}

En LoRa, la información puede existir físicamente pero no muestrearse con densidad suficiente. Para un estado aproximadamente constante y muestras condicionalmente independientes en $\mathcal T_{\rm obs}$,
\begin{equation}
\mat F_{\rm eff}=
\sum_{n\in\mathcal T_{\rm obs}}\mat F_{q,n}.
\end{equation}
Para respiración dinámica se empleará, en cambio, un modelo de estado y una FIM Bayesiana o posterior que incorpore correlación temporal. SenLoRa introduce métricas de disponibilidad en tráfico ambiente y estrategias para aumentar los símbolos útiles sin abandonar comunicación \citep{SenLoRa2025}. La tesis separará así \emph{información física por muestra}, \emph{cadencia disponible} e \emph{información acumulada}.

\section{Incertidumbre y abstención}

Se distinguirán el ruido aleatorio, la incertidumbre epistémica del estimador y la discrepancia del gemelo. La FIM caracteriza un límite local bajo un modelo supuesto; no cubre por sí sola el error de escena ni el cambio de dominio. El sistema informará una estimación, un intervalo calibrado y un índice de observabilidad, y podrá abstenerse cuando la evidencia sea insuficiente. La calibración se verificará mediante cobertura empírica y curvas de riesgo frente a cobertura, especialmente antes de interpretar una salida fisiológica.

\section{Información de la forma de onda y de la frecuencia respiratoria}

Para $x(t)=A\sin(2\pi f_rt+\varphi)$, la sensibilidad respecto de la tasa se
obtiene combinando la derivada del canal con
\begin{equation}
\frac{\partial x(t)}{\partial f_r}=2\pi At\cos(2\pi f_rt+\varphi).
\end{equation}
La información de $f_r$ depende de amplitud, duración, tiempos observados y
confusión con $A$, $\varphi$ y perturbaciones instrumentales. Si $A=0$, la
frecuencia no es identificable en este modelo, aunque $\partial H/\partial x$
sea grande. Por ello un mapa de sensibilidad del canal no basta para asegurar
una estimación de respiraciones por minuto.

En la validación externa, los índices predichos por el gemelo se contrastarán
con errores y fallos de estimadores sobre mediciones reservadas. Se distinguirá
el índice calculado con estado verdadero, utilizado como referencia explicativa,
del disponible con posición u orientación inciertas en inferencia.

\section{Efecto de eliminar perturbaciones sobre la sensibilidad útil}

Tras blanquear y apilar partes real e imaginaria de los jacobianos, sean
$A_q=\sqrt{2}[\Re(\Sigma_w^{-1/2}J_q)^T,
\Im(\Sigma_w^{-1/2}J_q)^T]^T$ y $A_\zeta$ su equivalente para las
perturbaciones. En el modelo local de covarianza fija,
\begin{equation}
F_q^{\rm eq}=A_q^T(I-A_\zeta A_\zeta^\dagger)A_q.
\end{equation}
La información útil corresponde a la sensibilidad que no puede explicarse
mediante esas perturbaciones. En particular, una fase instrumental arbitraria
por instante puede confundirse con el desplazamiento de un único camino.
La calibración basal o una referencia independiente permiten estudiar otra
condición de identificabilidad. Esta formulación fundamenta la prueba S6 de
preservación de movimiento después del procesamiento instrumental.


\section{Invariancia de fase y pérdida de sensibilidad física}
\label{sec:invariance_sensing_limit}

Para el observable espacial de la \cref{eq:spatial_gram_invariant}, una
perturbación pequeña satisface
\begin{equation}
\delta\mat G=(\delta\vect h)\vect h^H+\vect h(\delta\vect h)^H
+\mathcal O(\|\delta\vect h\|^2).
\end{equation}
Si $\delta\vect h=j\alpha\vect h$ con $\alpha\in\mathbb R$, los términos
lineales se cancelan. Para un cambio finito puramente rotacional la invariancia
es exacta. Por tanto, una representación que elimina una fase instrumental
común también puede eliminar una variación física que actúe en esa dirección.
Esto ocurre, por ejemplo, en la idealización de un único camino cuyo movimiento
sólo modifica su fase común y mantiene amplitud y vector espacial.

S4.7a contrastará fidelidad estática en el espacio invariante; S6 deberá evaluar
además $\partial\mat G/\partial x_H$ o la derivada del observable elegido,
comparándola con una referencia mecánica. La información de Fisher del observable
transformado debe derivarse de su distribución: los productos de canales
ruidosos no conservan automáticamente el modelo gaussiano aditivo ni su
covarianza original. Una buena concordancia de $G$ no garantiza que permanezca
información suficiente para recuperar respiración. Esta limitación conecta
la evaluación de invariantes con la proyección de perturbaciones y justifica
mantener una referencia de fase estable cuando la tarea física la requiera.

\section{Contraste experimental de la información preservada}
\label{sec:observable_jacobian_test}

La observabilidad se contrastará después de aplicar el operador realmente utilizado.
Para $\mathcal V$ no holomorfa, como magnitud o productos conjugados, se apilarán
partes reales e imaginarias del canal en $\mathbf y_{\mathbb R}$. Entonces
\begin{equation}
J_{\mathcal V,q}=\frac{\partial\mathcal V}{\partial\mathbf y_{\mathbb R}}
\frac{\partial\mathbf y_{\mathbb R}}{\partial\mathbf q},\qquad
\widehat J_{\mathcal V,q}^{(r)}=
\frac{\mathcal V(Y(\mathbf q+h\mathbf e_r))-
\mathcal V(Y(\mathbf q-h\mathbf e_r))}{2h}.
\label{eq:observable_real_jacobian}
\end{equation}
Se probarán varios pasos $h$ resolubles por la referencia mecánica para separar
error de truncamiento, histéresis y ruido amplificado por la diferencia finita.
La distancia cociente $\mathsf V_4$ exige analizar su geometría local o un representante
con referencia declarada; no se tratará como un vector de características arbitrario.

Se compararán dirección, norma y valores singulares del Jacobiano medido y simulado,
además de error relativo con un piso de denominador declarado. Si la sensibilidad
medida no se distingue del ruido, un error relativo pequeño no certifica un gemelo
útil. La FIM utilizará la covarianza del observable transformado y perturbaciones
instrumentales residuales; no se trasladará sin cambios el supuesto gaussiano del
CSI a productos o cocientes. Se confrontará su predicción con errores empíricos,
cobertura de intervalos y tasa de abstención en datos reservados.

Para dos personas, se estudiará el Jacobiano conjunto respecto de ambas fuentes y
sus direcciones casi colineales. Una tasa respiratoria correcta del componente
más fuerte no acredita separación de usuarios. Deben evaluarse asociación, errores
por persona y condiciones no identificables, antes de optimizar una red según esa
información. Los umbrales de confianza y de abstención se fijarán con validación.

La observabilidad establece qué puede recuperarse; no determina por sí sola cómo aprender
una representación robusta a partir de datos heterogéneos. El capítulo siguiente completa
el marco conceptual con aprendizaje autosupervisado, supervisión multimodal y adaptación
desde la simulación hacia las mediciones reales.
